\documentclass[11pt,a4paper]{article}

\usepackage[utf8]{inputenc}
\usepackage[T1]{fontenc}
\frenchspacing
\usepackage{lmodern}
\usepackage{geometry}
\geometry{top=2.5cm, bottom=2.5cm, left=2.5cm, right=2.5cm}

\usepackage{amsmath, amssymb}
\usepackage{booktabs}
\usepackage{float}
\usepackage{xcolor}
\usepackage{listings}
\usepackage{hyperref}
\usepackage{graphicx}
\usepackage{fancyhdr}
\usepackage{tikz}
\usetikzlibrary{shapes.multipart, arrows.meta, positioning, calc}

% Couleurs du rapport
\definecolor{MundiBlue}{RGB}{20, 45, 110}
\definecolor{LightGrey}{RGB}{245, 245, 247}
\definecolor{AlertRed}{RGB}{200, 30, 30}
\definecolor{CodeGreen}{RGB}{40, 140, 60}
\definecolor{CodeGray}{RGB}{100, 100, 100}

\hypersetup{
    colorlinks=true,
    linkcolor=MundiBlue,
    urlcolor=MundiBlue,
    citecolor=MundiBlue
}

% Configuration Listings pour Python
\lstset{
    language=Python,
    basicstyle=\ttfamily\small,
    backgroundcolor=\color{LightGrey},
    keywordstyle=\color{MundiBlue}\bfseries,
    stringstyle=\color{CodeGreen},
    commentstyle=\color{CodeGray}\itshape,
    numbers=left,
    numberstyle=\tiny\color{CodeGray},
    stepnumber=1,
    numbersep=8pt,
    showstringspaces=false,
    breaklines=true,
    frame=single,
    rulecolor=\color{black!15},
    tabsize=4
}

% Configuration des en-tetes et pieds de page (style TP-01)
\pagestyle{fancy}
\fancyhf{}
\fancyhead[L]{\small Université Mundiapolis}
\fancyhead[R]{\small Génie Informatique}
\fancyfoot[L]{\small Fellah Youssef}
\fancyfoot[R]{\small \thepage}
\renewcommand{\headrulewidth}{0.4pt}

\begin{document}

% ==========================================
% PAGE DE GARDE (STYLE OFFICIEL TP-01)
% ==========================================
\begin{titlepage}
    \thispagestyle{empty}
    \begin{center}
        \IfFileExists{logo_mundiapolis.png}{%
            \includegraphics[width=6.5cm]{logo_mundiapolis.png}\\[0.6cm]
        }{%
            \vspace*{2cm}
        }
    \end{center}

    \noindent
    \begin{minipage}[t]{0.5\textwidth}
        \textbf{UNIVERSITÉ MUNDIAPOLIS}\\
        École d’Ingénieurs\\
        Casablanca
    \end{minipage}
    \hfill
    \begin{minipage}[t]{0.5\textwidth}
        \begin{flushright}
            \textbf{2\textsuperscript{ème} Année CI}\\
            Génie Informatique\\
            Filière : Informatique
        \end{flushright}
    \end{minipage}

    \vspace{2cm}

    \noindent
    \rule{\textwidth}{1.5pt}\\[0.4cm]
    \begin{center}
        {\Huge \textbf{Principes de la Blockchain}}\\[0.4cm]
        {\Large Application des principes de la blockchain en Python}
    \end{center}
    \vspace{-0.2cm}
    \noindent
    \rule{\textwidth}{1.5pt}\\[1.5cm]

    \begin{center}
        {\large \textbf{Module : Blockchain --- TP 02}}\\[0.3cm]
        {\large \textbf{Preuve d'enjeu (Proof of Stake)}}
    \end{center}

    \vspace{3cm}

    \noindent
    \begin{minipage}[t]{0.5\textwidth}
        \textbf{Réalisé par :}\\
        \textsc{Fellah} Youssef
    \end{minipage}
    \hfill
    \begin{minipage}[t]{0.5\textwidth}
        \begin{flushright}
            \textbf{Encadré par :}\\
            Prof. \textsc{Louja} Ayoub
        \end{flushright}
    \end{minipage}

    \vfill

    \begin{center}
        \textbf{Année universitaire : 2026/2027}
    \end{center}
\end{titlepage}

% ==========================================
% TABLE DES MATIERES
% ==========================================
\newpage
\pagenumbering{roman}
\tableofcontents
\newpage
\pagenumbering{arabic}

% ==========================================
% CONTENU DU RAPPORT
% ==========================================

\section{Introduction}

\subsection{Contexte}
Dans les protocoles de consensus décentralisés, la \textbf{preuve d'enjeu} (\emph{Proof of Stake} ou PoS) constitue une alternative économe en énergie à la preuve de travail (\emph{Proof of Work}). Contrairement au PoW où la production d'un bloc résulte d'une compétition de calcul brute (minage par force brute du nonce), le protocole PoS sélectionne de manière pseudo-aléatoire le producteur de bloc (appelé \textbf{proposant} ou \emph{forger}).

La probabilité $p_i$ pour un validateur $i$ d'être désigné pour forger le prochain bloc est rigoureusement proportionnelle à la \textbf{mise} (\emph{stake}) qu'il a engagée et verrouillée en garantie dans le réseau :
\[
p_i = \frac{m_i}{\sum_{j} m_j} = \frac{m_i}{M}
\]
où $m_i$ désigne la mise du validateur $i$ et $M$ la masse totale des enjeux engagés par l'ensemble des validateurs actifs. 

Ce rapport présente l'implémentation, l'expérimentation et l'analyse de sécurité d'un réseau Proof of Stake complet développé en Python avec une architecture client-serveur HTTP, réalisé dans le cadre du TP~02 du module Blockchain.

\subsection{Objectifs}
Les travaux menés visent à :
\begin{enumerate}
    \item implémenter les règles d'adhésion au réseau avec seuil d'enjeu minimal (\textbf{Règle R1}) ;
    \item concevoir l'algorithme déterministe de sélection pondérée par la roulette des mises (\textbf{Règle R2}) ;
    \item orchestrer la production de blocs transparents et auditables par n'importe quel nœud (\textbf{Règles R3 et R4}) ;
    \item évaluer l'impact cumulatif des récompenses de forgeage et quantifier le risque de centralisation (\textbf{Règle R5}) ;
    \item auditer l'intégrité de la blockchain et identifier les blocs falsifiés de manière autonome (\textbf{Règle R6}) ;
    \item appliquer les sanctions économiques de confiscation de mise et d'exclusion (\emph{slashing}, \textbf{Règle R7}) ;
    \item tester la robustesse statistique du protocole face à des tricheurs exploitant des tirages biaisés.
\end{enumerate}

\subsection{Outils et Architecture}
Le réseau fonctionne selon une architecture distribuée interconnectée par le protocole HTTP :
\begin{itemize}
    \item \textbf{Serveur de consensus} (\texttt{Server.py}) : hébergé sur un nœud central sous le framework Flask (port 5000), maintenant l'état global du réseau, la liste des validateurs et le grand livre (\emph{ledger}) des blocs produits ;
    \item \textbf{Nœuds validateurs} (\texttt{main.ipynb} et notebooks individuels sous \texttt{participants/}) : clients autonomes utilisant la bibliothèque \texttt{requests} pour interagir avec l'API REST du serveur et réexécuter localement les règles de contrôle ;
    \item \textbf{Reproductibilité} : initialisation du générateur pseudo-aléatoire avec une graine fixe (\texttt{GRAINE = 2}) permettant la comparaison formelle des mesures empiriques avec les valeurs théoriques du protocole.
\end{itemize}

\subsection{Configuration initiale des validateurs}
Le réseau démarre avec quatre participants dont les mises initiales vérifient la contrainte $m_i \ge 10\text{ ETH}$ :

\begin{table}[H]
\centering
\begin{tabular}{@{}clcc@{}}
\toprule
\textbf{Nœud (NUM)} & \textbf{Validateur} & \textbf{Mise initiale ($m_i$)} & \textbf{Part théorique ($p_i$)} \\
\midrule
1 & Idriss  & 40 ETH & 40,0\,\% \\
2 & Moussa  & 25 ETH & 25,0\,\% \\
3 & Youssef & 20 ETH & 20,0\,\% \\
4 & Ayoub   & 15 ETH & 15,0\,\% \\
\midrule
\textbf{Total} & \textbf{Réseau ($M$)} & \textbf{100 ETH} & \textbf{100,0\,\%} \\
\bottomrule
\end{tabular}
\caption{État initial du pool de validateurs}
\label{tab:validateurs_init}
\end{table}

% ------------------------------------------------------------------
\section{Mise en place du protocole et Règles de Consensus}

\subsection{Adhésion et seuil d'engagement (Règle R1)}

\subsubsection*{Principe}
L'entrée d'un validateur requiert le séquestre d'une garantie minimale $\ge 10\text{ ETH}$. Cette garantie sert de collatéral économique : elle permet de sanctionner financièrement tout comportement byzantin (malveillant). Le point de terminaison \texttt{POST /rejoindre} vérifie cette contrainte et met à jour la mise totale du réseau.

\subsubsection*{Résultats}
Les quatre participants enregistrent successivement leurs enjeux :
\begin{lstlisting}
200 {'mise_totale': 40, 'validateur': {'blocs': 0, 'exclu': False, 'mise': 40, 'nom': 'Idriss'}}
200 {'mise_totale': 65, 'validateur': {'blocs': 0, 'exclu': False, 'mise': 25, 'nom': 'Moussa'}}
200 {'mise_totale': 85, 'validateur': {'blocs': 0, 'exclu': False, 'mise': 20, 'nom': 'Youssef'}}
200 {'mise_totale': 100, 'validateur': {'blocs': 0, 'exclu': False, 'mise': 15, 'nom': 'Ayoub'}}
\end{lstlisting}
Une tentative d'adhésion avec une mise de 5~ETH est immédiatement rejetée par le serveur :
\begin{center}
\texttt{400 \{'erreur': 'mise inferieure a 10'\}}
\end{center}

% ------------------------------------------------------------------
\subsection{Sélection du proposant par roulette (Règle R2)}

\subsubsection*{Principe de la roulette déroulée}
Avant la création de chaque bloc, le protocole effectue un tirage aléatoire $t \in [0 ; M]$. On modélise le capital total comme un segment continu $[0 ; M]$ découpé en intervalles successifs dont la largeur correspond exactement à la mise de chaque validateur. Le validateur sélectionné est celui dont l'intervalle cumulé contient $t$.

\begin{figure}[H]
\centering
\begin{tikzpicture}[
    scale=0.11,
    segment/.style={rectangle, draw=MundiBlue, thick, minimum height=1.3cm, font=\sffamily\scriptsize, align=center},
    axe/.style={-{Stealth[length=2.5mm]}, thick, MundiBlue}
]
\node[segment, fill=LightGrey, minimum width=40cm] (idriss) at (20, 0) {Idriss (40 ETH)\\$[0 ; 40]$};
\node[segment, fill=LightGrey!70, minimum width=25cm] (moussa) at (52.5, 0) {Moussa (25 ETH)\\$]40 ; 65]$};
\node[segment, fill=LightGrey!50, minimum width=20cm] (youssef) at (75, 0) {Youssef (20 ETH)\\$]65 ; 85]$};
\node[segment, fill=LightGrey!30, minimum width=15cm] (ayoub) at (92.5, 0) {Ayoub (15 ETH)\\$]85 ; 100]$};

% Axe cumulé
\draw[axe] (0, -2.5) -- (105, -2.5) node[right, font=\scriptsize\sffamily]{Mise cumulée};
\foreach \x/\label in {0/0, 40/40, 65/65, 85/85, 100/100}
    \draw[thick] (\x, -2.0) -- (\x, -3.0) node[below, font=\tiny\sffamily]{\label};

% Tirage test 52.3
\draw[AlertRed, thick, -{Stealth[length=2mm]}] (52.3, 3.2) -- (52.3, 1.2) node[above=0.6cm, font=\scriptsize\bfseries]{$t = 52{,}3$};
\node[AlertRed, font=\tiny\sffamily] at (52.3, -4.5) {$40 < 52{,}3 \le 65 \implies$ Moussa};
\end{tikzpicture}
\caption{Mécanisme de sélection par roulette pondérée (projection linéaire)}
\label{fig:roulette}
\end{figure}

\subsubsection*{Algorithme}
La fonction \texttt{choisir\_proposant} est une fonction pure : pour un ensemble ordonné de mises et un tirage donné, elle produit invariablement le même résultat, permettant son exécution locale chez n'importe quel nœud vérificateur.

\begin{lstlisting}
def choisir_proposant(mises, tirage):
    cumul = 0.0
    for nom, mise in mises:
        cumul += mise
        if tirage <= cumul:
            return nom
    return mises[-1][0] if mises else None
\end{lstlisting}

\subsubsection*{Validation sur tirages de référence}
L'évaluation sur les valeurs tests prévues par le sujet confirme l'exactitude de l'intervalle :
\begin{itemize}
    \item $t = 52{,}3 \in ]40 ; 65] \implies$ \textbf{Moussa}
    \item $t = 12{,}0 \in [0 ; 40] \implies$ \textbf{Idriss}
    \item $t = 95{,}5 \in ]85 ; 100] \implies$ \textbf{Ayoub}
\end{itemize}

% ------------------------------------------------------------------
\subsection{Production d'un bloc et auditabilité (Règles R3 et R4)}

\subsubsection*{Structure d'un bloc PoS}
Contrairement au bloc PoW qui stocke un nonce miné, le bloc PoS consigne la preuve de sa légitimité en intégrant explicitement :
\begin{itemize}
    \item \texttt{numero} : hauteur du bloc dans la chaîne ;
    \item \texttt{proposant} : identité du validateur autorisé à signer ;
    \item \texttt{tirage} : valeur aléatoire générée par le consensus ;
    \item \texttt{mises} : cliché (\emph{snapshot}) complet des mises actives au moment du tirage.
\end{itemize}

\subsubsection*{Vérification autonome}
L'appel à \texttt{POST /produire\_bloc} génère le premier bloc :
\begin{lstlisting}
{
  'numero': 1,
  'proposant': 'Ayoub',
  'tirage': 95.6,
  'mises': [['Idriss', 40], ['Moussa', 25], ['Youssef', 20], ['Ayoub', 15]]
}
\end{lstlisting}
Tout nœud recalcule localement $\texttt{choisir\_proposant}(\text{mises}, 95{,}6) = \text{"Ayoub"}$, attestant que le bloc respecte la règle R3.

\subsubsection*{Distribution sur un échantillon court (15 blocs)}
La production de 14 blocs supplémentaires permet d'obtenir les 15 premiers blocs de la chaîne (tableau~\ref{tab:15_blocs}).

\begin{table}[H]
\centering
\small
\begin{tabular}{@{}clrcl@{}}
\toprule
\textbf{Bloc} & \textbf{Proposant} & \textbf{Tirage ($t$)} & \textbf{Intervalle associé} & \textbf{Statut R4} \\
\midrule
1  & Ayoub   & 95,60 & $]85 ; 100]$ & Conforme \\
2  & Ayoub   & 94,78 & $]85 ; 100]$ & Conforme \\
3  & Idriss  &  5,66 & $[0 ; 40]$   & Conforme \\
4  & Idriss  &  8,49 & $[0 ; 40]$   & Conforme \\
5  & Youssef & 83,55 & $]65 ; 85]$  & Conforme \\
6  & Youssef & 73,60 & $]65 ; 85]$  & Conforme \\
7  & Youssef & 66,97 & $]65 ; 85]$  & Conforme \\
8  & Idriss  & 30,81 & $[0 ; 40]$   & Conforme \\
9  & Moussa  & 60,59 & $]40 ; 65]$  & Conforme \\
10 & Moussa  & 60,68 & $]40 ; 65]$  & Conforme \\
11 & Moussa  & 58,12 & $]40 ; 65]$  & Conforme \\
12 & Idriss  & 15,84 & $[0 ; 40]$   & Conforme \\
13 & Moussa  & 43,07 & $]40 ; 65]$  & Conforme \\
14 & Idriss  & 39,35 & $[0 ; 40]$   & Conforme \\
15 & Youssef & 72,30 & $]65 ; 85]$  & Conforme \\
\bottomrule
\end{tabular}
\caption{Historique des 15 premiers blocs générés ($\text{seed} = 2$)}
\label{tab:15_blocs}
\end{table}

Sur 15 blocs, la distribution observée est : Idriss (5 blocs, 33,3\,\%), Moussa (4 blocs, 26,7\,\%), Youssef (4 blocs, 26,7\,\%) et Ayoub (2 blocs, 13,3\,\%). Sur un échantillon aussi restreint, la variance statistique est forte : Moussa et Youssef obtiennent le même nombre de blocs alors que leurs mises diffèrent de 5~ETH.

% ------------------------------------------------------------------
\subsection{Convergence statistique sur 5 000 blocs (Règle R2)}

Pour valider expérimentalement la convergence de la loi des grands nombres ($p_i \to m_i/M$), 5\,000 blocs additionnels sont produits, portant la chaîne à 5\,015 blocs sans récompense active ($r = 0$).

\begin{table}[H]
\centering
\begin{tabular}{@{}lcccc@{}}
\toprule
\textbf{Validateur} & \textbf{Mise} & \textbf{Part théorique ($p_i$)} & \textbf{Part observée} & \textbf{Écart absolu} \\
\midrule
Idriss  & 40 ETH & 40,0\,\% & 39,5\,\% & $-0{,}5\,\%$ \\
Moussa  & 25 ETH & 25,0\,\% & 25,2\,\% & $+0{,}2\,\%$ \\
Youssef & 20 ETH & 20,0\,\% & 20,7\,\% & $+0{,}7\,\%$ \\
Ayoub   & 15 ETH & 15,0\,\% & 14,6\,\% & $-0{,}4\,\%$ \\
\bottomrule
\end{tabular}
\caption{Comparaison des fréquences de production sur 5\,015 blocs}
\label{tab:5000_blocs}
\end{table}

L'écart maximal observé n'excède pas $0{,}7$ point de pourcentage. La proportionnalité du tirage est donc expérimentalement vérifiée.

% ------------------------------------------------------------------
\subsection{Dynamique des récompenses et risque de centralisation (Règle R5)}

\subsubsection*{Formulation mathématique}
La règle R5 alloue une récompense au validateur qui forge un bloc : sa mise augmente d'une fraction $r$ :
\[
m_i^{(k+1)} = m_i^{(k)} \cdot (1 + r)
\]
tandis que les mises des autres validateurs restent stationnaires.

\subsubsection*{Résultats sur 1\,000 blocs ($r = 0{,}3\,\%$)}
L'expérience est lancée sur un réseau réinitialisé avec les mêmes 4 validateurs sur 1\,000 blocs avec $r = 0{,}003$.

\begin{table}[H]
\centering
\small
\begin{tabular}{@{}lrrrrcc@{}}
\toprule
\textbf{Validateur} & \textbf{Mise dép.} & \textbf{Mise fin} & \textbf{Part dép.} & \textbf{Part fin} & \textbf{Blocs} & \textbf{Facteur ($\times$)} \\
\midrule
Idriss  &  40,0 ETH & 175,2 ETH & 40,0\,\% & \textbf{62,9\,\%} & 493 & \textbf{4,38} \\
Moussa  &  25,0 ETH &  50,2 ETH & 25,0\,\% & 18,0\,\%          & 233 & 2,01 \\
Youssef &  20,0 ETH &  31,3 ETH & 20,0\,\% & 11,2\,\%          & 149 & 1,56 \\
Ayoub   &  15,0 ETH &  21,8 ETH & 15,0\,\% &  7,8\,\%          & 125 & 1,45 \\
\midrule
\textbf{Réseau} & \textbf{100,0 ETH} & \textbf{278,5 ETH} & \textbf{100,0\,\%} & \textbf{100,0\,\%} & \textbf{1000} & \textbf{2,78} \\
\bottomrule
\end{tabular}
\caption{Évolution des mises et des parts après 1\,000 blocs avec récompenses ($r = 0{,}3\,\%$)}
\label{tab:recompenses}
\end{table}

\subsubsection*{Analyse de l'amplification de centralisation}
Le capital global du réseau a cru d'un facteur 2,78. Bien qu'aucun validateur n'ait perdu d'ETH dans l'absolu (Moussa a par exemple doublé son capital, passant de 25 à 50,2~ETH), la part relative de Moussa, Youssef et Ayoub s'est effondrée. 

Seul Idriss progresse plus vite que la moyenne du réseau ($\times 4{,}38 > 2{,}78$), portant sa part de 40\,\% à près de 63\,\%. C'est l'illustration empirique du biais de centralisation \emph{"rich get richer"} inhérent au Proof of Stake naïf : détenir une mise plus forte confère plus de blocs, donc plus de récompenses, amplifiant l'avantage initial à chaque génération.

% ------------------------------------------------------------------
\subsection{Détection des fraudes et Slashing (Règles R6 et R7)}

\subsubsection*{Attaque par usurpation d'identité (Étape 8)}
Un validateur malveillant (Ayoub) force l'inscription d'un bloc à son propre nom via \texttt{POST /falsifier}, alors que le tirage attribué ($t = 73{,}60$) se situe dans l'intervalle de Youssef ($]65 ; 85]$). Le bloc falsifié porte le numéro 6.

\subsubsection*{Algorithme de vérification autonome (Règle R6)}
Tout participant audite l'ensemble de la chaîne localement sans solliciter l'avis du serveur :

\begin{lstlisting}
def verifier_bloc(bloc):
    # R6 : recalcule le proposant attendu d'apres le tirage public
    attendu = choisir_proposant(bloc["mises"], bloc["tirage"])
    return bloc["proposant"] == attendu

def verifier_chaine(chaine):
    return [b["numero"] for b in chaine if not verifier_bloc(b)]
\end{lstlisting}

L'exécution locale produit le verdict détaillé (tableau~\ref{tab:audit_fraude}).

\begin{table}[H]
\centering
\small
\begin{tabular}{@{}clrcl@{}}
\toprule
\textbf{Bloc} & \textbf{Signé par} & \textbf{Tirage} & \textbf{Recalcul local R2} & \textbf{Verdict d'audit} \\
\midrule
1 & Ayoub   & 95,60 & Ayoub   & \texttt{ok} \\
2 & Ayoub   & 94,78 & Ayoub   & \texttt{ok} \\
3 & Idriss  &  5,66 & Idriss  & \texttt{ok} \\
4 & Idriss  &  8,49 & Idriss  & \texttt{ok} \\
5 & Youssef & 83,55 & Youssef & \texttt{ok} \\
\textbf{6} & \textbf{Ayoub} & \textbf{73,60} & \textbf{Youssef} & \textbf{\texttt{FALSIFIE}} \\
7 & Youssef & 66,97 & Youssef & \texttt{ok} \\
8 & Idriss  & 30,81 & Idriss  & \texttt{ok} \\
9 & Moussa  & 60,59 & Moussa  & \texttt{ok} \\
\bottomrule
\end{tabular}
\caption{Audit autonome de la chaîne par recalcul public}
\label{tab:audit_fraude}
\end{table}

Le bloc 6 est immédiatement et formellement identifié comme falsifié :
\begin{center}
\texttt{recalcul local : [6] \quad vs \quad réponse du réseau : [6]}
\end{center}

\subsubsection*{Exécution de la sanction (Règle R7)}
Dès la détection, la sanction est appliquée via \texttt{POST /sanctionner} avec $\text{fraction} = 0{,}5$ :
\begin{lstlisting}
Ayoub : {'exclu': True, 'mise_restante': 7.5, 'perte': 7.5}
mise totale apres exclusion : 85.0 ETH
\end{lstlisting}
Ayoub subit une pénalité financière immédiate de $7{,}5\text{ ETH}$ (50\,\% de sa mise) et est banni de la liste active des validateurs (\texttt{exclu = True}). La masse totale d'enjeux retombe à 85~ETH, augmentant automatiquement la part de marché des trois validateurs honnêtes restants.

% ------------------------------------------------------------------
\section{Tests approfondis du réseau}

\subsection{Test 1 : Précision selon la taille de l'échantillon}
La simulation est répétée sur 500 blocs puis sur 100\,000 blocs avec les 4 validateurs initiaux :

\begin{table}[H]
\centering
\begin{tabular}{@{}lcccc@{}}
\toprule
\textbf{Validateur} & \textbf{Part théorique} & \textbf{500 blocs} & \textbf{100\,000 blocs} & \textbf{Écart (100k)} \\
\midrule
Idriss  & 40,0\,\% & 39,4\,\% & 39,9\,\% & $-0{,}1\,\%$ \\
Moussa  & 25,0\,\% & 24,6\,\% & 25,0\,\% & $\pm 0{,}0\,\%$ \\
Youssef & 20,0\,\% & 20,0\,\% & 20,1\,\% & $+0{,}1\,\%$ \\
Ayoub   & 15,0\,\% & 16,0\,\% & 15,0\,\% & $\pm 0{,}0\,\%$ \\
\bottomrule
\end{tabular}
\caption{Convergence de la distribution en fonction du nombre de blocs}
\label{tab:convergence_test}
\end{table}

L'écart maximal passe de $1{,}0$ point sur 500 blocs à $0{,}1$ point sur 100\,000 blocs, démontrant la stricte conformité asymptotique de l'algorithme.

\subsection{Test 2 : Arrivée de 3 nouveaux validateurs}
Trois validateurs rejoignent le réseau : Aymen (100~ETH), Timan (50~ETH) et Meryem (50~ETH). La mise totale passe à 300~ETH. Sur 5\,000 blocs :
\begin{itemize}
    \item Aymen détient $100/300 = 33{,}3\,\%$ de la mise et forge $33{,}4\,\%$ des blocs ;
    \item Timan et Meryem détiennent chacun $16{,}7\,\%$ et forgent respectivement $17{,}3\,\%$ et $16{,}3\,\%$ ;
    \item Les 4 participants d'origine conservent leur mise absolue mais voient leur part relative divisée exactement par 3.
\end{itemize}

\subsection{Test 3 : Emballement par récompense accrue ($r = 1{,}0\,\%$)}
En portant le taux de récompense de $0{,}3\,\%$ à $1{,}0\,\%$ sur 1\,000 blocs, l'effet d'amplification devient exponentiel :
\begin{itemize}
    \item Idriss atteint une mise finale de \textbf{93\,914,3 ETH} (pour 40 ETH au départ), capturant \textbf{99,9\,\%} de l'ensemble de la masse monétaire du réseau ;
    \item Moussa passe à 69,0~ETH, Youssef à 37,1~ETH et Ayoub à 26,2~ETH : leur capital s'est apprécié mais leur part relative devient nulle ($0{,}0\,\%$).
\end{itemize}
Ce test démontre qu'un taux de rémunération excessif détruit la décentralisation en quelques centaines de blocs.

\subsection{Test 4 : Attaque par tirage biaisé et détection fréquentielle}

\subsubsection*{Scénario d'attaque}
Ayoub n'usurpe plus le nom d'autrui sur un tirage public défavorable. À la place, il forge lui-même des tirages artificiels situés systématiquement dans son propre intervalle $[85 ; 100]$ (1 bloc sur 4).

\subsubsection*{Incapacité du contrôle R6 unitaire}
Chaque bloc étant individuellement cohérent avec le tirage qu'il transporte ($\texttt{verifier\_bloc} = \text{True}$), la liste des blocs falsifiés reste totalement vide :
\begin{center}
\texttt{rep["blocs\_falsifies"] = []}
\end{center}

\subsubsection*{Détection par analyse statistique}
L'attaque est démasquée par l'analyse fréquentielle globale (règle $\text{observe} - \text{attendu} > 3{,}0\,\%$) :

\begin{table}[H]
\centering
\begin{tabular}{@{}lcccc@{}}
\toprule
\textbf{Validateur} & \textbf{Part de mise} & \textbf{Part des blocs} & \textbf{Écart} & \textbf{Verdict de surveillance} \\
\midrule
Idriss  & 40,0\,\% & 28,8\,\% & $-11{,}2\,\%$ & normal \\
Moussa  & 25,0\,\% & 18,8\,\% & $-6{,}2\,\%$  & normal \\
Youssef & 20,0\,\% & 16,2\,\% & $-3{,}8\,\%$  & normal \\
\textbf{Ayoub} & \textbf{15,0\,\%} & \textbf{36,2\,\%} & $\mathbf{+21{,}2\,\%}$ & \textbf{\texttt{SUSPECT}} \\
\bottomrule
\end{tabular}
\caption{Détection de l'attaque discrète par anomalie de distribution}
\label{tab:frequences_suspect}
\end{table}

Alors qu'Ayoub ne dispose que de 15\,\% de la mise, il génère plus de 36\,\% des blocs. Les deux modes de contrôle sont strictement complémentaires : la vérification unitaire intercepte les falsifications syntaxiques, tandis que le moniteur de fréquence neutralise la manipulation de l'aléatoire.

% ------------------------------------------------------------------
\section{Synthèse}

\subsection{Bilan}
Ce travail de simulation a permis de concrétiser les concepts fondamentaux du Proof of Stake :
\begin{itemize}
    \item la \textbf{sélection proportionnelle par roulette}, garantissant l'équité mathématique sans gaspillage d'énergie ;
    \item la \textbf{vérification décentralisée R6}, rendue possible par la publicité des tirages et des états de mise enregistrés dans chaque bloc (R4) ;
    \item l'\textbf{effet de centralisation économique R5}, inhérent aux mécanismes de récompense proportionnels ;
    \item le \textbf{slashing R7}, assurant l'alignement des incitations en pénalisant financièrement la trahison.
\end{itemize}

\subsection{Limites de la simulation}
Le réseau implémenté présente des simplifications pédagogiques :
\begin{itemize}
    \item \textbf{génération de l'aléatoire} : le tirage est confié au producteur ou au serveur HTTP, ce qui permet la fraude du test 4. En environnement réel, l'aléa doit provenir d'une \textbf{fonction aléatoire vérifiable} (VRF) ou d'un schéma de \emph{commit-reveal} (ex. RANDAO sur Ethereum) ;
    \item \textbf{signatures cryptographiques} : les identités reposent sur des chaînes de caractères au lieu de clés publiques et signatures asymétriques (ECDSA / BLS) ;
    \item \textbf{topologie centralisée} : les nœuds communiquent avec un serveur Flask unique au lieu d'un réseau peer-to-peer (P2P) à diffusion par commérage (\emph{gossip protocol}).
\end{itemize}

\subsection{Améliorations possibles}
\begin{itemize}
    \item intégrer une fonction aléatoire vérifiable (VRF) cryptographique garantissant l'impossibilité de truquer le tirage ;
    \item implémenter un mécanisme de pondération non linéaire (ex. preuve d'enjeu quadratique ou délégation LPoS) pour freiner la centralisation des récompenses ;
    \item coupler les blocs avec des signatures numériques ECDSA et des arbres de Merkle pour les transactions.
\end{itemize}

% ------------------------------------------------------------------
\section{Annexes \& Bibliographie}

\subsection{Annexe A : Dépôt du projet}

Le code source intégral (notebooks exécutés, implémentation des règles et lanceur serveur) est consultable sur le dépôt GitHub :
\begin{center}
\url{https://github.com/ind4x/BLOCKCHAIN-LAB/tree/main/TP-02}
\end{center}

\subsection{Annexe B : Glossaire}
\begin{description}
    \item[Proof of Stake (PoS)] Protocole de consensus attribuant la production de blocs selon la proportion de cryptomonnaie mise en gage.
    \item[Mise (Stake)] Montant d'actifs numériques immobilisé par un validateur pour garantir sa probité sous peine de confiscation.
    \item[Proposant (Forger)] Validateur tiré au sort pour construire, signer et diffuser le prochain bloc.
    \item[Slashing] Mécanisme de pénalité confisquant une fraction ou la totalité de la mise d'un nœud en cas de faute avérée.
    \item[Roulette des mises] Méthode d'échantillonnage découpant un segment continu proportionnellement aux enjeux individuels.
    \item[VRF (Verifiable Random Function)] Primitive cryptographique générant un nombre aléatoire imprévisible assorti d'une preuve publique de validité.
\end{description}

\subsection{Bibliographie}
\begin{thebibliography}{9}
\bibitem{king} S.~King et S.~Nadal, \emph{PPCoin: Peer-to-Peer Crypto-Currency with Proof-of-Stake}, 2012.
\bibitem{buterin} V.~Buterin et V.~Griffith, \emph{Casper the Friendly Finality Gadget}, arXiv:1710.09437, 2017.
\bibitem{micali} Y.~Gilad, R.~Hemo, S.~Micali, G.~Vlachos et N.~Zeldovich, \emph{Algorand: Scaling Byzantine Agreements for Cryptocurrencies}, ACM SOSP, 2017.
\bibitem{flask} Pallets Projects, \emph{Flask Documentation (3.1.x) --- A lightweight WSGI web application framework}, 2026.
\bibitem{tp} Université Mundiapolis, \emph{TP 02 --- Preuve d'enjeu (Proof of Stake)}, Campus Nouaceur, Filière Ingénierie Informatique.
\end{thebibliography}

\end{document}
